\chapter{Архитектура контура управления}

Контур управления строится на двух архитектурных решениях: его компоненты
размещаются вне федерации, а взаимодействуют друг с другом только через
каталог кластеров. Первое решение делает проверяемым сценарий отказа
кластера, второе~--- обеспечивает честное сравнение планировщиков.

\section{Звенья контура}

Контур образует цепочку
\begin{center}
наблюдаемость $\to$ SLO $\to$ каталог кластеров $\to$ планировщик $\to$
контроллер аварийного восстановления,
\end{center}
в которой каждое звено служит источником данных для следующего. Метрики
кластеров агрегируются и превращаются в показатели уровня обслуживания;
каталог хранит атрибуты кластеров и их текущее состояние; планировщик
выбирает кластер для новой нагрузки по данным каталога; контроллер
аварийного восстановления обнаруживает отказ кластера, переключает трафик
и сообщает об отказе каталогу, откуда эту информацию получает
планировщик.

\begin{figure}[ht]
  \centering
  \begin{tikzpicture}[
      font=\scriptsize, >={Stealth[length=2mm]},
      comp/.style={draw, rounded corners=2pt, fill=white, align=center,
                   minimum height=8mm, inner sep=2pt, text width=24mm},
      grp/.style={draw, rounded corners=4pt, dashed, inner sep=5pt},
      lbl/.style={font=\scriptsize\itshape, align=center, inner sep=1pt}]
    % управляющий кластер
    \node[comp] (gw) {Шлюз\\Envoy Gateway};
    \node[comp, left=16mm of gw] (dr) {Контроллер\\восстановления};
    \node[comp, right=16mm of gw] (q) {Thanos Querier\\и Ruler};
    \node[comp, below=10mm of gw] (reg) {Каталог\\FastAPI + PostgreSQL};
    \node[comp, below=10mm of dr] (sch) {Планировщик\\meta | naive};
    \node[comp, below=10mm of q] (gr) {Grafana, Jaeger,\\Prometheus};
    \node[comp, above=8mm of gw] (k6) {k6, генератор\\нагрузки (хост)};
    \begin{scope}[on background layer]
      \node[grp, fit=(dr)(q)(sch)(gr), label={[font=\scriptsize\bfseries]above left:управляющий кластер \texttt{control}}] (ctl) {};
    \end{scope}
    \draw[->] (k6) -- (gw);
    \draw[->] (dr) -- (gw) node[lbl, midway, above] {веса};
    \draw[->] (dr) -- (reg) node[lbl, pos=0.45, below left] {\texttt{down}};
    \draw[->] (q) -- (reg) node[lbl, pos=0.45, below right] {SLI};
    \draw[<->] (sch) -- (reg) node[lbl, midway, below=1mm] {кандидаты};
    % кластеры федерации
    \foreach \c/\x/\s in {a/-30mm/left, b/30mm/right} {
      \node[comp] (w\c) at ($(reg)+(\x,-30mm)$) {\texttt{app}, поток \texttt{w-<n>}\\frontend $\to$ backend};
      \node[comp, below=2mm of w\c] (m\c) {Prometheus + sidecar,\\OTel Collector};
      \node[comp, below=2mm of m\c] (p\c) {Kyverno, Chaos Mesh};
      \begin{scope}[on background layer]
        \node[grp, fit=(w\c)(m\c)(p\c), label={[font=\scriptsize\bfseries]below:\texttt{cluster-\c}}] (cl\c) {};
      \end{scope}
      \coordinate (d\c) at ($(cl\c.north)+(-4mm,0)$);
      \coordinate (u\c) at ($(cl\c.north)+(4mm,0)$);
      \draw[->, thick] (d\c |- ctl.south) -- (d\c);
      \draw[<-, thick] (u\c |- ctl.south) -- (u\c);
    }
    \node[lbl, left=5mm] at ($(cla.north)!0.5!(cla.north |- ctl.south)$) {трафик,\\размещение};
    \node[lbl, right=5mm] at ($(cla.north)!0.5!(cla.north |- ctl.south)$) {метрики,\\пробы};
    \node[lbl, left=5mm] at ($(clb.north)!0.5!(clb.north |- ctl.south)$) {трафик,\\размещение};
    \node[lbl, right=5mm] at ($(clb.north)!0.5!(clb.north |- ctl.south)$) {метрики,\\пробы};
  \end{tikzpicture}
  \caption{Архитектура контура управления на стенде}
  \label{fig:architecture}

  \todo{рисунок требует проверки}
\end{figure}

Управляющий контур размещается вне федерации, в отдельном управляющем
кластере (рис.~\ref{fig:architecture}). Это существенно: компонент,
размещённый внутри одного из управляемых кластеров, перестаёт работать
именно при отказе этого кластера, и сценарий отказа становится
непроверяемым. Внутри кластеров федерации находятся только нагрузка,
локальный сбор метрик и трейсов, движок политик и средство внесения
отказов.

\section{Роли компонентов}

Компоненты контура делятся на три роли (табл.~\ref{tab:roles}).

\begin{table}[ht]
  \centering
  \caption{Роли компонентов контура}
  \label{tab:roles}
  \small
  \begin{tabularx}{\textwidth}{@{}lXX@{}}
    \toprule
    Роль & Компоненты & Назначение \\
    \midrule
    Источники истины & Prometheus и Thanos, каталог кластеров &
      Что происходит с сервисом; какие кластеры существуют и в каком они
      состоянии \\
    Решающие & Планировщик, контроллер аварийного восстановления &
      Куда разместить нагрузку; объявлять ли инцидент. Решения
      записываются в журнал \\
    Исполнители & Kyverno, шлюз Envoy Gateway &
      Пропустить или отклонить размещение; куда направить трафик. Решений
      не принимают \\
    \bottomrule
  \end{tabularx}
\end{table}

\section{Каталог как общая точка взаимодействия}

Контроллер аварийного восстановления сообщает об отказе кластера не
планировщику напрямую, а каталогу, присваивая кластеру статус
\texttt{down}. Планировщик читает каталог перед каждым решением. Таким
образом компоненты не знают друг о друге, и второго источника истины о
состоянии кластера не возникает. Фоновый опрос метрик каталогом не
затирает статус \texttt{down}: его устанавливает и снимает только
контроллер аварийного восстановления.

Фильтрация кластеров по обязательным ограничениям также выполняется
каталогом, и ей пользуются оба планировщика~--- и разрабатываемый, и
наивный. Поэтому разница между планировщиками состоит строго в учёте
динамических метрик.

\section{Инварианты}

Нарушение следующих свойств считается дефектом, а не компромиссом:
\begin{enumerate}
  \item Наблюдатель находится вне наблюдаемого: ни один компонент контура не
        размещается в кластере федерации.
  \item Каталог~--- единственный источник состояния для планировщика; к
        системе мониторинга напрямую планировщик не обращается.
  \item Оба планировщика соблюдают обязательные ограничения; различаются
        они только учётом предпочтений.
  \item Два рубежа контроля политик независимы: при отключённом
        планировщике политики внутри кластеров продолжают действовать.
  \item Каждое решение о размещении записывается со всеми кандидатами и их
        оценками, а не только с выбранным кластером.
  \item Отсутствие метрик (\texttt{unknown}) не равно отказу кластера
        (\texttt{down}) и не является основанием для переключения трафика.
  \item Время внесения отказа берётся из инструмента внесения отказа, а не
        из наблюдений контура; в противном случае время обнаружения отказа
        было бы тождественно нулевым.
\end{enumerate}

\section{Схема данных каталога}

Каталог хранит данные в пяти таблицах (табл.~\ref{tab:schema}).

\begin{table}[ht]
  \centering
  \caption{Таблицы каталога кластеров}
  \label{tab:schema}
  \small
  \begin{tabularx}{\textwidth}{@{}lXX@{}}
    \toprule
    Таблица & Содержимое & Назначение \\
    \midrule
    \texttt{clusters} & Регион, провайдер, сертификации, метки, статус &
      Инвентарь и состояние кластеров \\
    \texttt{cluster\_metrics} & Снимок доли ошибок, задержки p95, загрузки и
      свободных ресурсов раз в 30~с & Входные данные планировщика \\
    \texttt{placement\_decisions} & Все кандидаты с оценками, снимком метрик
      и причиной отсева; тип планировщика; исход & Анализ причин выбора и
      офлайн-пересчёт при других параметрах \\
    \texttt{policy\_violations} & Нарушения политик, источник обнаружения,
      допустимые кластеры & Журнал нарушений \\
    \texttt{incidents} & Фазы цикла восстановления и время внесения отказа &
      Расчёт времени обнаружения и восстановления \\
    \bottomrule
  \end{tabularx}
\end{table}

Существенно, что в журнале решений хранится полный список кандидатов со
снимком их метрик на момент решения. Без этого сравнительный анализ свёлся бы
к сравнению итоговых показателей без понимания того, почему планировщик
выбрал именно так, а анализ чувствительности к параметрам потребовал бы
повторения экспериментов.
